The \code{Thread} class represents a single thread of execution
in the target process. Any Process has at least one Thread, and
multi-threaded target processes may have more. Each Thread has an associated
integral value known as its LWP, which serves as a handle for communicating
with the OS about the thread (e.g., a PID value on Linux). On some systems,
depending on availability, a Thread may have information from the user space
threading library.

\begin{apient}
Thread::ptr
Thread::const_ptr
\end{apient}

\apidesc{
The Thread::ptr and Thread::const\_ptr respectively represent a pointer and a
const pointer to a Thread object. Both pointer types are reference counted
and cause the underlying Thread object to be cleaned when there are no more
references. ProcControlAPI maintains internal references to any Thread it
actively controls, relinquishing those references when the thread exits or is
detached.
}

\begin{apient}
Dyninst::LWP getLWP() const
\end{apient}

\apidesc{
This function returns an OS handle for this thread. On Linux this returns a
pid\_t for this thread. On FreeBSD, this returns a lwpid\_t.
}

\begin{apient}
Process::ptr getProcess()
Process::const_ptr getProcess() const
\end{apient}

\apidesc{
These functions return a pointer to the Process object that contains this
thread. 
}

\begin{apient}
bool isStopped() const
\end{apient}

\apidesc{
This function returns true if this thread is in a stopped state and false
otherwise.
}

\begin{apient}
bool isRunning() const
\end{apient}

\apidesc{
This function returns true if this thread is in a running state and false
otherwise.
}

\begin{apient}
bool isLive() const
\end{apient}

\apidesc{
This function returns true if this thread is alive, and it returns false if this
thread has been destroyed.
}

\begin{apient}
bool isDetached() const
\end{apient}

\apidesc{
This function returns true if this thread has been detached via
Process::temporaryDetach and false otherwise.
}

\begin{apient}
bool isInitialThread() const
\end{apient}

\apidesc{
This function returns true if this thread is the initial thread for the process
and false otherwise.
}

\begin{apient}
bool stopThread()
\end{apient}

\apidesc{
This function moves the thread to into a stopped state. Upon a successful call
to this function the Thread object will be paused and will not resume execution
until the Thread is continued. It is an error to call this function from a
callback. Instead of calling this function, a callback can stop a thread by
returning Process::cbThreadStop or Process::cbProcStop.
ProcControlAPI may deliver callbacks when this function is called.
Upon success this function returns true, otherwise it returns false. Upon an
error a subsequent call to getLastError returns details on the error.
}

\begin{apient}
bool continueThread()
\end{apient}

\apidesc{
This function moves the thread into a running state. It is an error to call
this function from a callback. Instead of calling this function, a callback
can stop a thread by returning Process::cbThreadContinue or
Process::cbProcContinue.
ProcControlAPI may deliver callbacks when this function is called.
Upon success this function returns true, otherwise it returns false. Upon an
error a subsequent call to getLastError returns details on the error.
}

\begin{apient}
bool getRegister(Dyninst::MachRegister reg, Dyninst::MachRegisterVal &val) const
\end{apient}

\apidesc{
This function gets the value of a single register from this thread. The
register is specified by the reg parameter, and the value of the register is
returned by the val parameter.

It is an error to call this function on a thread that is not in the stopped
state.

Upon success this function returns true, otherwise it returns false. Upon an
error a subsequent call to getLastError returns details on the error.
}

\begin{apient}
bool getAllRegisters(RegisterPool pool) const
\end{apient}

\apidesc{
This function reads the values of every register in the thread and returns them
as part of the RegisterPool object pool. Depending on the OS, this call may
be more efficient that calling Thread::getRegister multiple times.
It is an error to call this function on a thread that is not in the stopped
state.
Upon success this function returns true, otherwise it returns false. Upon an
error a subsequent call to getLastError returns details on the error.
}

\begin{apient}
bool setRegister(Dyninst::MachRegister reg, Dyninst::MachRegisterVal val) const
\end{apient}

\apidesc{
This function writes the value of a single register in this thread. The
register is specified by the reg parameter, and the value that should be
written is specified by the val parameter.

It is an error to call this function on a thread that is not in the stopped
state.
Upon success this function returns true, otherwise it returns false. Upon an
error a subsequent call to getLastError returns details on the error.
}

\begin{apient}
bool setAllRegisters(RegisterPool &pool) const
\end{apient}

\apidesc{
This function sets the values of every register in this thread to the values
specified in the RegisterPool object pool. Depending on the OS, this call may
be more efficient that calling Thread::setRegister multiple times.
It is an error to call this function on a thread that is not in the stopped
state.
Upon success this function returns true, otherwise it returns false. Upon an
error a subsequent call to getLastError returns details on the error.
}

\begin{apient}
bool haveUserThreadInfo() const
\end{apient}

\apidesc{
This function returns true if information about this
Thread\'s underlying user-level thread is available.
}

\begin{apient}
Dyninst::THR_ID getTID() const
\end{apient}

\apidesc{
This function returns the unique identifier for the user-level thread. This
value is only valid if haveUserThreadInfo returns true.
}

\begin{apient}
Dyninst::Address getStartFunction() const
\end{apient}

\apidesc{
This function returns the address of the initial function for the user-level
thread. This value is only valid if haveUserThreadInfo returns true.
}

\begin{apient}
Dyninst::Address getStackBase() const
\end{apient}

\apidesc{
This function returns the address of the bottom of the user-level
thread\'s stack. This value is only valid if
haveUserThreadInfo returns true.
}

\begin{apient}
unsigned long getStackSize() const
\end{apient}

\apidesc{
This function returns the size in bytes of the user-level
thread\'s allocated stack. This value is only valid if
haveUserThreadInfo returns true.
}

\begin{apient}
Dyninst::Address getTLS() const
\end{apient}

\apidesc{
This function returns the address of the user-level
thread\'s thread local storage area. This value is only
valid if haveUserThreadInfo returns true.
}

\begin{apient}
bool readThreadLocalMemory(void *buffer,
                           Library::const_ptr lib,
                           Dyninst::Offset tls_symbol_offset,
                           size_t size) const
\end{apient}

\apidesc{
This function reads from a symbol in thread local storage (TLS) memory. TLS is
memory that is local to a thread and has a lifetime matching the thread. The
tls\_symbol\_offset is the TLS symbol\'s offset in lib, and
can be found by reading a TLS symbol\'s value. The lib
parameter can point to a library or the executable. The buffer parameter
specifies an address in the controller process where ProcControlAPI should
write the copied bytes. 
It is an error to call this function on a Thread that is not in a stopped state.
It is also an error to call this function on a Thread that has not have
user-level thread information, which can be tested with haveUserThreadInfo.
This function returns true on success and false on error. Upon an error a
subsequent call to getLastError returns details on the error.
}

\begin{apient}
bool writeThreadLocalMemory(Library::const_ptr lib,
                            Dyninst::Offset tls_symbol_offset,
                            const void *buffer,
                            size_t size) const
\end{apient}

\apidesc{
This function writes to a symbol in thread local storage (TLS) memory. TLS is
memory that is local to a thread and has a lifetime matching the thread. The
tls\_symbol\_offset is the TLS symbol\'s offset in lib, and
can be found by reading a TLS symbol\'s value. The lib
parameter can point to a library or the executable. The buffer parameter
specifies an address in the controller process where ProcControlAPI should read
the bytes to be copied.
It is an error to call this function on a Thread that is not in a stopped state.
It is also an error to call this function on a Thread that has not have
user-level thread information, which can be tested with haveUserThreadInfo.
This function returns true on success and false on error. Upon an error a
subsequent call to getLastError returns details on the error.
}

\begin{apient}
bool getThreadLocalAddress(Library::const_ptr lib,
                           Dyninst::Offset tls_symbol_offset,
                           Dyninst::Address &result_addr) const
\end{apient}

\apidesc{
This function looks up the address of a symbol in thread local storage (TLS)
memory. The tls\_symbol\_offset is the TLS symbol\'s offset
in lib, and can be found by reading a TLS symbol\'s value.
The lib parameter can point to a library or the executable. The result\_addr
parameter will be set to the target address for the TLS symbol in this Thread.
It is an error to call this function on a Thread that is not in a stopped state.
It is also an error to call this function on a Thread that has not have
user-level thread information, which can be tested with haveUserThreadInfo.
This function returns true on success and false on error. Upon an error a
subsequent call to getLastError returns details on the error.
}

\begin{apient}
bool postIRPC(IRPC::ptr irpc) const
\end{apient}

\apidesc{
This function posts the given irpc to the Thread. The IRPC is put irpc into
the Thread\'s queue of posted IRPCs and will be run when
ready.
Each instance of an IRPC object can be posted at most once. It is an error to
attempt to post a single IRPC object multiple times.
This function returns true on success and false on error. Upon an error a
subsequent call to getLastError returns details on the error.
}

\begin{apient}
bool getPostedIRPCs(std::vector<IRPC::ptr> &rpcs) const
\end{apient}

\apidesc{
This function returns all IRPCs posted to this Thread in the vector rpcs. This
does not include any running IRPC.
This function returns true on success and false on error. Upon an error a
subsequent call to getLastError returns details on the error.
}

\begin{apient}
IRPC::const_ptr getRunningIRPC() const
\end{apient}

\apidesc{
This function returns a const pointer to any IRPC that is actively running on
this Thread. If there is no IRPC actively running, then this function returns
IRPC::const\_ptr().
}

\begin{apient}
void setSingleStepMode(bool mode) const
\end{apient}

\apidesc{
This function sets whether a Thread is in single-step mode. If called with a
mode of true, then the Thread is put in single-step mode. If called with a
mode of false, then the Thread is taken out of single-step mode. 
A Thread in single-step mode will pause execution at each instruction and
trigger an EventSingleStep event. After each EventSingleStep is handled (and
presuming the Thread is still running and in single-step mode) the Thread will
execute one more instruction and trigger another EventSingleStep.
}

\begin{apient}
bool getSingleStepMode() const
\end{apient}

\apidesc{
This function returns true if the Thread is in single-step mode and false
otherwise.
}

\begin{apient}
void *getData() const
void setData(void *p)
\end{apient}

\apidesc{
These functions respectively get and set an opaque data object that can be
associated with this Thread. The data is not interpreted by ProcControlAPI,
but remains associated with the Thread.
}
